% ===========================================================
% 3. PRELIMINARIES
% ===========================================================
\section{Preliminaries}
\label{sec:preliminaries}

Versatile embodied navigation requires a single agent to fulfill
heterogeneous user goals---ranging from precise metric coordinates to
abstract semantic descriptions---using only on-board visual observations
and producing executable motion commands at the control rate. Before
describing ABot-N1, we cast this requirement as a unified
conditional sequential decision problem, specialize the goal
specification to the five tasks considered in this work, and introduce
the VLM-conditioned action-policy abstraction that subsequent sections
build on.

\subsection{Embodied Navigation As Goal-Conditioned Visual Control}
\label{sec:prelim-formulation}

We model an embodied navigator as a discrete-time agent acting in an
unknown environment $\mathcal{E}$. At each step $t$, the agent receives
an ego-centric RGB observation
$I_t \in \mathbb{R}^{N\times H\times W\times 3}$ from $N$ cameras and an
optional proprioceptive/odometry signal $q_t$, and is conditioned on a
task specification $g$ that encodes the user's intent. Following recent
navigation foundation
models~\cite{zhang2025navfom,zhang2024uninavid,physicalintelligence2024pi05},
we formulate the navigator as a goal-conditioned visual policy
\begin{equation}
\pi_\theta:\;\bigl(g,\,I_{1:t},\,q_{1:t}\bigr)\;\longmapsto\; a_{t:t+H}
\label{eq:nav-policy}
\end{equation}
which predicts a chunk of $H$ future low-level commands
$a_{t:t+H}=(a_t,a_{t+1},\ldots,a_{t+H-1})$ from the visual history
$I_{1:t}$ and the goal condition $g$. Each command
\begin{equation}
a_i=\bigl(x_i,\,y_i,\,\sin\theta_i,\,\cos\theta_i,\,c_i\bigr)\in\mathbb{R}^4\times\{0,1\}
\label{eq:waypoint}
\end{equation}
is a continuous $\mathrm{SE}(2)$ waypoint in the agent's local frame
augmented with a binary completion flag $c_i$ that triggers the
``arrive'' signal once the goal is reached. The agent rolls out
$\pi_\theta$ in a closed loop, re-planning at every step from updated
observations, until $c_i=1$ or a step budget is exhausted; an episode
succeeds when the final pose lies within a task-specific tolerance of
the true target. Crucially, this formulation is map-free: $\pi_\theta$
is given neither a global metric map nor depth or semantic
segmentation, and must extract all necessary geometry and semantics
from the on-the-fly RGB stream alone, matching the constraints of
real deployment on commodity
sensors~\cite{zhang2024navid,zhang2024uninavid}.

\subsection{Unifying FIVE Navigation Tasks in ONE Framework}
\label{sec:prelim-tasks}

The goal condition $g$ is realized as a natural-language instruction
$g=\ell$, where $\ell$ is a free-form sentence that encodes the
navigation intent. This single interface absorbs
all five tasks studied in this report without per-task heads or
bespoke encoders: in \textit{point-goal}, $\ell$ encodes a relative
metric offset; in \textit{instruction-following}
(R2R/RxR/VLN-CE-style episodes), $\ell$ is a step-by-step natural-language
route description; in \textit{object-goal}, $\ell$ is an open-vocabulary
category phrase such as ``find the nearest chair''; in \textit{POI-goal},
$\ell$ names a point of interest characterized by a landmark or
business identity; and in \textit{person-following}, $\ell$ describes
the target person via attribute or identity clauses.
Casting heterogeneous goal modalities into a unified language instruction
$\ell$ lets a single policy $\pi_\theta$ share parameters and
representations across tasks~\cite{zhang2025navfom,abotn0,zhang2024uninavid},
which is a prerequisite for transferring semantic priors learned on one
task to another.

\subsection{VLM-Conditioned Policy and Continuous-Action Decoding}
\label{sec:prelim-vlm}
Realizing $\pi_\theta$ at the scale required for open-world
generalization motivates initializing the policy from a pre-trained
vision-language
model~\cite{bai2023qwenvl,chen2024spatialvlm,physicalintelligence2024pi05}.
A natural design factorizes the navigator into a \emph{VLM reasoner}
$f_\phi$ and a lightweight \emph{action policy} $\psi_\xi$:
\begin{align}
s_t &= f_\phi\!\bigl(I_{1:t},\,\ell\bigr),
\label{eq:vlm-reason}\\
a_{t:t+H} &= \psi_\xi\!\bigl(s_t,\,I_{1:t},\ell,\,q_t\bigr)
\label{eq:action-head}
\end{align}
where $s_t$ denotes an intermediate guidance signal produced by the
reasoner--its concrete form (latent features, decoded language, visual
anchors, or a combination thereof) is a design choice deferred to
specific instantiations.
This factorization brings two benefits.
First, the broad linguistic and visual priors acquired during
web-scale VLM pre-training transfer directly into navigation,
providing open-vocabulary recognition, spatial-relation understanding,
and instruction parsing without task-specific
modules~\cite{bai2023qwenvl,chen2024spatialvlm,cheng2024spatialrgpt}.
Second, decoupling the heavy reasoner from the reactive action policy
allows them to operate at \emph{asymmetric computational
budgets}---a principle that ABot-N1 materializes as a slow-fast
dual-system architecture detailed in \S\ref{sec:method}.
